\newpage
\section{Implementacja}

Implementacja projektu \texttt{PZ\_Projekt\_Listy} została wykonana w języku C++. Program wykorzystuje listę dwukierunkową, której elementy są przechowywane w pamięci dynamicznej. Każdy element zawiera wartość liczbową oraz wskaźniki umożliwiające przechodzenie do następnego i poprzedniego elementu.

Kod projektu został podzielony na pliki nagłówkowe \texttt{.hpp} oraz pliki źródłowe \texttt{.cpp}. W przedstawionej implementacji wykorzystano klasy \texttt{Node} i \texttt{List}, a także odwołania do klas \texttt{NodeFactory} oraz \texttt{Iterator}.

\subsection{Implementacja klasy Node}

Klasa \texttt{Node} reprezentuje pojedynczy element listy. Przechowuje wartość liczbową w polu \texttt{data} oraz dwa wskaźniki: \texttt{next} i \texttt{previous}. Pierwszy wskazuje następny element listy, a drugi element poprzedni.

Deklaracja klasy znajduje się w pliku \texttt{Node.hpp}.

\begin{lstlisting}[language=C++, caption={Deklaracja klasy Node}]
#ifndef Node_hpp
#define Node_hpp

#include <stdio.h>

class Node{
public:
    int data;
    Node *next;
    Node *previous;

    Node(int number);
};

#endif /* Node_hpp */
\end{lstlisting}

Konstruktor klasy \texttt{Node} przyjmuje wartość liczbową, którą zapisuje w polu \texttt{data}. Oba wskaźniki są inicjalizowane wartością \texttt{nullptr}, ponieważ nowo utworzony element nie jest jeszcze połączony z innymi elementami listy.

\begin{lstlisting}[language=C++, caption={Konstruktor klasy Node}]
#include "Node.hpp"

Node::Node(int number){
    data = number;
    next = nullptr;
    previous = nullptr;
}
\end{lstlisting}

Takie rozwiązanie pozwala na późniejsze połączenie elementu z pozostałymi węzłami listy poprzez ustawienie odpowiednich wskaźników.

\subsection{Inicjalizacja i zwalnianie pamięci}

Konstruktor klasy \texttt{List} ustawia wskaźniki \texttt{first} oraz \texttt{last} na \texttt{nullptr}. Oznacza to, że lista początkowo nie zawiera żadnych elementów. Zmienna \texttt{size} otrzymuje wartość zero.

\begin{lstlisting}[language=C++, caption={Konstruktor i destruktor klasy List}]
List::List(){
    first = nullptr;
    last = nullptr;
    size = 0;
}

List::~List(){
    clear();
}
\end{lstlisting}

Destruktor wywołuje metodę \texttt{clear()}, która usuwa wszystkie elementy listy i zwalnia zajmowaną przez nie pamięć. Dzięki temu po zakończeniu działania obiektu wykonywane jest czyszczenie dynamicznie zaalokowanych węzłów.

Do sprawdzania, czy lista jest pusta, służy metoda \texttt{isEmpty()}.

\begin{lstlisting}[language=C++, caption={Sprawdzanie, czy lista jest pusta}]
bool List::isEmpty(){
    return first == nullptr;
}
\end{lstlisting}

Metoda zwraca wartość \texttt{true}, gdy wskaźnik \texttt{first} nie wskazuje na żaden element. W przeciwnym przypadku zwracana jest wartość \texttt{false}.

\subsection{Dodawanie elementów do listy}

Metoda \texttt{addFirst()} dodaje nowy element na początku listy. W tym celu tworzy węzeł za pomocą metody \texttt{NodeFactory::create()}. Jeżeli lista jest pusta, nowy element staje się jednocześnie pierwszym i ostatnim elementem. W przeciwnym przypadku zostaje połączony z dotychczasowym początkiem listy.

\begin{lstlisting}[language=C++, caption={Dodawanie elementu na początku listy}]
void List::addFirst(int value){
    Node *node = NodeFactory::create(value);

    if (isEmpty()){
        first = last = node;
    }else{
        node -> next = first;
        first->previous = node;
        first = node;
    }

    size++;
    iterator.setCurrent(first);
}
\end{lstlisting}

Po dodaniu elementu zwiększana jest wartość \texttt{size}. Następnie aktualizowana jest pozycja iteratora za pomocą metody \texttt{setCurrent()}.

Dodawanie elementu na końcu listy realizuje metoda \texttt{addLast()}. Nowy węzeł zostaje połączony z dotychczasowym ostatnim elementem. Wskaźnik \texttt{last} zostaje następnie ustawiony na nowy węzeł.

\begin{lstlisting}[language=C++, caption={Dodawanie elementu na końcu listy}]
void List::addLast(int value){
    Node *node = NodeFactory::create(value);

    if(isEmpty()){
        first = last = node;
    }else{
        node->previous = last;
        last->next = node;
        last = node;
    }

    size++;
    iterator.setCurrent(first);
}
\end{lstlisting}

Obie metody obsługują przypadek pustej listy, dzięki czemu pierwszy dodawany element prawidłowo ustawia oba wskaźniki graniczne.

Metoda \texttt{addAt()} umożliwia wstawienie elementu pod wskazanym indeksem. Najpierw sprawdzana jest poprawność indeksu. Wartości mniejsze od zera oraz większe od aktualnego rozmiaru listy są odrzucane. Dla indeksu zerowego wywoływana jest metoda \texttt{addFirst()}, natomiast dla indeksu równego rozmiarowi listy — metoda \texttt{addLast()}.

\begin{lstlisting}[language=C++, caption={Sprawdzanie indeksu w metodzie addAt}]
void List::addAt(int index, int value){
    if(index < 0 || index > size){
        cout << "Nieprawidlowy indeks!" << endl;
        return;
    }

    if(index == 0){
        addFirst(value);
        return;
    }

    if(index == size){
        addLast(value);
        return;
    }
\end{lstlisting}

W przypadku wstawiania elementu wewnątrz listy program przechodzi do węzła znajdującego się pod wskazanym indeksem. Następnie tworzy nowy element i aktualizuje połączenia między sąsiednimi węzłami.

\begin{lstlisting}[language=C++, caption={Wstawianie elementu wewnątrz listy}]
    Node *current = first;

    for(int i = 0; i < index; i++){
        current = current->next;
    }

    Node *node = NodeFactory::create(value);

    node->next = current;
    node->previous = current->previous;
    current->previous->next = node;
    current->previous = node;

    size++;
    iterator.setCurrent(first);
}
\end{lstlisting}

Dzięki aktualizacji obu kierunków połączeń lista zachowuje swoją dwukierunkową strukturę.

\subsection{Wyświetlanie elementów listy}

Do wyświetlania zawartości listy od początku do końca służy metoda \texttt{showAll()}. Przechodzi ona przez kolejne węzły, zaczynając od \texttt{first}, aż do momentu napotkania wartości \texttt{nullptr}. Każdy element jest wyświetlany wraz z numerem porządkowym.

\begin{lstlisting}[language=C++, caption={Wyświetlanie listy w przód}]
void List::showAll(){
    int counter = 1;
    Node *current = first;

    while(current != nullptr){
        cout << counter << ") "
             << current->data << endl;

        current = current->next;
        counter++;
    }
}
\end{lstlisting}

Wyświetlanie w odwrotnym kierunku realizuje metoda \texttt{showReverse()}. W tym przypadku przechodzenie rozpoczyna się od wskaźnika \texttt{last}, a kolejne elementy są odczytywane za pomocą wskaźnika \texttt{previous}.

\begin{lstlisting}[language=C++, caption={Wyświetlanie listy w odwrotnej kolejności}]
void List::showReverse(){
    int counter = 1;
    Node *current = last;

    while(current != nullptr){
        cout << counter << ") "
             << current->data << endl;

        current = current->previous;
        counter++;
    }
}
\end{lstlisting}

Wykorzystanie obu wskaźników pozwala na przechodzenie przez tę samą listę w dwóch kierunkach bez konieczności tworzenia jej dodatkowej kopii.

\subsection{Usuwanie elementów}

Metoda \texttt{removeFirst()} usuwa pierwszy element listy. Na początku sprawdza, czy lista nie jest pusta. Jeśli nie zawiera elementów, wyświetlany jest odpowiedni komunikat i metoda kończy działanie.

\begin{lstlisting}[language=C++, caption={Usuwanie pierwszego elementu}]
void List::removeFirst(){
    if(isEmpty()){
        cout << "Lista jest pusta!" << endl;
        return;
    }

    Node *temp = first;

    if(first == last){
        first = nullptr;
        last = nullptr;
    }else{
        first = first->next;
        first->previous = nullptr;
    }

    delete temp;
    size--;
    iterator.setCurrent(first);
}
\end{lstlisting}

W przypadku listy zawierającej tylko jeden element oba wskaźniki zostają ustawione na \texttt{nullptr}. W przeciwnym przypadku początek listy zostaje przesunięty na kolejny węzeł, a jego wskaźnik \texttt{previous} zostaje wyzerowany. Usunięty węzeł jest zwalniany operatorem \texttt{delete}.

Metoda \texttt{removeLast()} wykonuje analogiczną operację na końcu listy.

\begin{lstlisting}[language=C++, caption={Usuwanie ostatniego elementu}]
void List::removeLast(){
    if(isEmpty()){
        cout << "Lista jest pusta!" << endl;
        return;
    }

    Node *temp = last;

    if(first == last){
        first = nullptr;
        last = nullptr;
    }else{
        last = last->previous;
        last->next = nullptr;
    }

    delete temp;
    size--;
    iterator.setCurrent(first);
}
\end{lstlisting}

W przypadku listy wieloelementowej wskaźnik \texttt{last} zostaje przesunięty na poprzedni węzeł. Następnie jego wskaźnik \texttt{next} jest ustawiany na \texttt{nullptr}.

Usuwanie elementu pod określonym indeksem realizuje metoda \texttt{removeAt()}. Najpierw sprawdzana jest poprawność indeksu. Jeżeli wskazany element znajduje się na początku lub na końcu listy, wykorzystywane są odpowiednio metody \texttt{removeFirst()} i \texttt{removeLast()}.

\begin{lstlisting}[language=C++, caption={Sprawdzanie indeksu w metodzie removeAt}]
void List::removeAt(int index){
    if(index < 0 || index >= size){
        cout << "Nieprawidlowy indeks!" << endl;
        return;
    }

    if(index == 0){
        removeFirst();
        return;
    }

    if(index == size - 1){
        removeLast();
        return;
    }
\end{lstlisting}

Dla elementu znajdującego się wewnątrz listy program odnajduje odpowiedni węzeł i łączy ze sobą jego poprzednika oraz następnika. Następnie usuwa wskazany element i aktualizuje rozmiar listy.

\begin{lstlisting}[language=C++, caption={Usuwanie elementu ze środka listy}]
    Node *current = first;

    for(int i = 0; i < index; i++){
        current = current->next;
    }

    current->previous->next = current->next;
    current->next->previous = current->previous;

    delete current;
    size--;
    iterator.setCurrent(first);
}
\end{lstlisting}

Dzięki odpowiedniemu przepięciu wskaźników usunięcie węzła nie powoduje przerwania połączeń pomiędzy pozostałymi elementami listy.

\subsection{Obsługa iteratora}

W klasie \texttt{List} wykorzystano obiekt \texttt{iterator}, który umożliwia przechodzenie do kolejnych i poprzednich elementów. Metody \texttt{showNext()} oraz \texttt{showPrevious()} sprawdzają, czy możliwe jest wykonanie danego kroku.

\begin{lstlisting}[language=C++, caption={Przejście do następnego elementu}]
void List::showNext(){
    if(iterator.hasNext()){
        cout << "Nastepny element: "
             << iterator.next() << endl;
    }else{
        cout << "Brak nastepnego elementu!"
             << endl;
    }
}
\end{lstlisting}

Jeżeli iterator może przejść do następnego elementu, metoda \texttt{next()} zwraca jego wartość. W przeciwnym przypadku wyświetlany jest komunikat o braku następnego elementu.

Analogicznie działa przechodzenie w przeciwnym kierunku.

\begin{lstlisting}[language=C++, caption={Przejście do poprzedniego elementu}]
void List::showPrevious(){
    if(iterator.hasPrevious()){
        cout << "Poprzedni element: "
             << iterator.previous() << endl;
    }else{
        cout << "Brak poprzedniego elementu!"
             << endl;
    }
}
\end{lstlisting}

Po wykonaniu operacji dodawania lub usuwania elementów pozycja iteratora jest ustawiana ponownie za pomocą \texttt{iterator.setCurrent(first)}. Pełna implementacja działania iteratora znajduje się w osobnych plikach projektu.

\subsection{Czyszczenie listy}

Metoda \texttt{clear()} usuwa wszystkie elementy listy. Przechodzi przez kolejne węzły, zapamiętuje adres aktualnego elementu, przesuwa wskaźnik na następny węzeł, a następnie usuwa zapamiętany element.

\begin{lstlisting}[language=C++, caption={Usuwanie wszystkich elementów listy}]
void List::clear(){
    Node *current = first;

    while(current != nullptr){
        Node *temp = current;
        current = current->next;
        delete temp;
    }

    first = nullptr;
    last = nullptr;
    size = 0;
    iterator.setCurrent(nullptr);
}
\end{lstlisting}

Po zakończeniu pętli wskaźniki \texttt{first} i \texttt{last} są ustawiane na \texttt{nullptr}, a rozmiar listy zostaje wyzerowany. Iterator również otrzymuje wartość \texttt{nullptr}, co odpowiada pustej liście.

\subsection{Podsumowanie implementacji}

Przedstawiona implementacja umożliwia dodawanie, usuwanie i wyświetlanie elementów listy dwukierunkowej. Wykorzystanie wskaźników \texttt{next} i \texttt{previous} umożliwia przechodzenie w obu kierunkach, natomiast wskaźniki \texttt{first} i \texttt{last} zapewniają bezpośredni dostęp do początku i końca struktury.

Tworzenie nowych węzłów odbywa się przez \texttt{NodeFactory::create()}, a obsługa przechodzenia pomiędzy elementami została powiązana z obiektem \texttt{iterator}. Pamięć zajmowana przez elementy jest zwalniana podczas ich usuwania oraz przy czyszczeniu całej listy.
